\documentclass[10pt,twocolumn]{article}
\usepackage[margin=1in]{geometry}
\usepackage{graphicx}
\usepackage{booktabs}
\usepackage{amsmath}
\usepackage{amssymb}
\usepackage{hyperref}
\usepackage{xcolor}
\usepackage{caption}
\usepackage{subcaption}
\usepackage{float}

\hypersetup{colorlinks=true, linkcolor=blue, citecolor=blue, urlcolor=blue}

\title{\textbf{012: Ternary Triadic Networks for\\Efficient Rotation-Robust Vision and Streaming Inference}}
\author{
  [Author Name]\\
  [Institution / Independent]\\
  \texttt{gbranaa4@gmail.com}
}
\date{\today}

\newcommand{\tritnet}{TritCognition}
\newcommand{\tritlm}{TritLM}

\begin{document}
\maketitle

\begin{abstract}
We present \textbf{012}, a neural computing stack grounded in ternary triadic
logic, where each unit operates across three streams ---
\textit{Observer} (0), \textit{Shadow} (1), and \textit{Light} (2) ---
corresponding to suppression, detection, and relational context.
We make five empirical contributions.
(1)~\textbf{Rotation robustness}: \tritnet{} achieves lower normalized
stability loss than parameter-matched float32 and ternary baselines on CIFAR-10
and STL-10, with a controlled ablation confirming triadic structure --- not
ternary quantization --- as the cause.
(2)~\textbf{Fixed memory}: \tritlm{}, a ternary triadic language model,
maintains a constant 0.05\,KB working state regardless of context length,
versus a GPT KV-cache that grows to 150\,MB at 100k tokens
($3.2\times10^6\times$ reduction).
(3)~\textbf{Noise correction}: clean-context trit memory recovers
$+0.34$ nats at 50\% character corruption, with recovery increasing
monotonically with corruption severity.
(4)~\textbf{Forgetting resistance}: streaming new content does not degrade
performance on prior content ($\Delta{=}{-}0.004$ nats, slight improvement).
(5)~\textbf{Interpretability}: discrete trit memory states produce
measurably distinct patterns per content type (stage directions vs.\ dialogue:
cosine similarity $0.14$), a property unavailable in float32 models.
All hardware is specified in synthesizable SystemVerilog with 36/36
test cases passing.
\end{abstract}

\section{Introduction}

Binary computing is approaching fundamental limits. MOSFET scaling has slowed,
power density has plateaued, and the energy cost of running large neural
networks is a growing industrial crisis~\cite{thompson2021}.
The theoretical optimal integer base for arithmetic is $e \approx 2.718$,
making base-3 the most information-dense practical integer
system~\cite{hayes2001}.

We propose \textbf{012} --- a complete computing stack built on ternary logic
and triadic neural architecture --- motivated by three converging observations:

\textbf{Physical.} Ternary transistors storing three stable voltage states have
been demonstrated at lab scale~\cite{shukla2019}.

\textbf{Mathematical.} Ternary weight networks $\{-1, 0, +1\}$ eliminate
multiplications, reducing energy per operation by an estimated
60--70\%~\cite{rastegari2016, andri2018}.

\textbf{Structural.} Biological neurons exhibit three functional modes ---
inhibition, baseline, excitation --- mapping naturally to the 012 trit.
Predictive coding theory~\cite{rao1999, friston2005} suggests the brain
computes prediction errors rather than binary activations.

\noindent Our contributions:
\begin{enumerate}
  \item \textbf{\tritnet}: triadic CNN with ternary weights, predictive
        coding loss, spatial attention, and working memory gating.
  \item \textbf{Ablation suite}: 7 models isolate triadic structure,
        ternary quantization, and each architectural component.
  \item \textbf{\tritlm}: ternary triadic language model with
        fixed-size recurrent state replacing the KV cache.
  \item \textbf{Five streaming proofs}: memory footprint, noise correction,
        forgetting resistance, and trit interpretability.
  \item \textbf{Hardware}: synthesizable SystemVerilog, 36/36 testbench pass.
\end{enumerate}

\section{Method}

\subsection{Ternary Weight Quantization}

\begin{equation}
  W_\text{trit} = \begin{cases}
    \text{sign}(W) & |W| > \tau \\
    0              & \text{otherwise}
  \end{cases}, \quad \tau = 0.7 \cdot \mathbb{E}[|W|]
\end{equation}

Gradients flow via straight-through estimator (STE)~\cite{bengio2013}:
$\partial\mathcal{L}/\partial W \approx \partial\mathcal{L}/\partial
W_\text{trit} \cdot \mathbf{1}[|W| \leq 1]$.
Training: float warmup (epochs 1--8), then ternary (epochs 9--50).
Adaptive threshold $\tau$ prevents network collapse; fixed thresholds
produce $>$95\% zero weights.

\subsection{Triadic Convolutional Block}

\begin{align}
  s_0 &= \sigma(\text{Conv}_{1\times1}(x)) \in [0,1]
         \;\; \text{(Observer)} \\
  s_1 &= \tanh(\text{Conv}_{3\times3}(x)) \in [-1,1]
         \;\; \text{(Shadow)} \\
  s_2 &= \tanh(\text{Conv}_{5\times5}(x)) \in [-1,1]
         \;\; \text{(Light)}
\end{align}
\begin{equation}
  \text{out} = s_1(1-s_0) + s_2 s_0
\end{equation}

The Observer gate continuously decides whether each spatial location
is understood locally ($s_1$, 3$\times$3 field) or relationally
($s_2$, 5$\times$5 field). This dynamic routing over scales
--- without any explicit symmetry constraint --- produces rotation
robustness as an emergent property.
Hardware consensus gate: $\text{consensus}(a,b,c) = \text{sign}(a+b+c)$.

\subsection{Predictive Coding Loss}

\begin{equation}
  \mathcal{L} = \mathcal{L}_\text{CE} +
  \lambda\sum_i\text{MSE}(\hat{x}_{i+1}, x_{i+1}),
  \quad \lambda=0.01
\end{equation}

Each triadic block predicts its own input via a 1$\times$1 head.
Ablation confirms $+1.54$pp normalized stability contribution on STL-10.

\subsection{TritMemoryCell}

\begin{align}
  f_t &= \sigma(W_f \cdot x_t) \quad \text{(forget gate)} \\
  w_t &= \tanh(W_w \cdot x_t) \quad \text{(write gate)} \\
  r_t &= \sigma(W_r \cdot x_t) \quad \text{(read gate)} \\
  M_t &= M_{t-1}(1-f_t) + w_t f_t \\
  \text{out}_t &= x_t + r_t \cdot M_t
\end{align}

All gate weights are ternary $\{-1,0,+1\}$. The cell state $M$ persists
across forward passes, enabling streaming inference with $O(1)$ memory.
Trit interpretation: $M_i{=}-1$ (pattern suppressed),
$M_i{=}0$ (unknown), $M_i{=}+1$ (pattern activated).

\subsection{\tritlm{} --- Triadic Attention}

Standard transformers use three float32 projections $Q,K,V$.
\tritlm{} replaces these with triadic projections:

\begin{align}
  s_0 &= \sigma(W_0 x) \quad \text{Observer: gates attention scores} \\
  s_1 &= \tanh(W_1 x) \quad \text{Shadow: content queries} \\
  s_2 &= \tanh(W_2 x) \quad \text{Light: relational keys+values}
\end{align}
\begin{align}
  \text{scores} &= (s_1 \cdot s_2^\top / \sqrt{d}) \cdot \bar{s}_0 \\
  \text{attn}   &= \text{softmax}(\text{scores} + \text{causal mask}) \\
  \text{out}    &= \text{attn} \cdot s_2
\end{align}

Each block also contains a TritMemoryCell that persists state across
tokens, replacing the KV cache entirely. Total inference memory:
$O(1)$ in sequence length.

\subsection{\tritnet{} Architecture}

Input $\to$ [TB(3$\to$32)$\to$Pool] $\to$ [TB(32$\to$64)$\to$Pool]
$\to$ [TB(64$\to$128)$\to$Pool] $\to$ SpatialAttn $\to$ GAP
$\to$ TritMemoryCell $\to$ Linear.
Total: \textbf{396,174} parameters vs.\ ResNet18: 11,181,642.

\section{Experiments}

\subsection{Setup}

\begin{table}[h]
\centering\small
\begin{tabular}{ll}
\toprule
Config & Value \\ \midrule
Datasets      & CIFAR-10, STL-10 \\
Train samples & 50k (CIFAR-10), 5k (STL-10) \\
Test samples  & 10k (CIFAR-10), 8k (STL-10) \\
Epochs        & 50 (warmup:8, ternary:42) \\
Optimizer     & SGD (ResNet18); Adam (others) \\
LR schedule   & CosineAnnealingLR \\
Rotation      & Bilinear, fill=0, angles $\{0,45,90,135,180,270\}^\circ$ \\
Device        & NVIDIA RTX 5060 (CUDA 13.2) \\
\bottomrule
\end{tabular}
\end{table}

\subsection{Rotation Robustness}

\begin{table}[H]
\centering\small
\caption{Accuracy (\%) and normalized stability on CIFAR-10 under rotation.
Stability $\downarrow$ = $(0^\circ - \text{worst}) / 0^\circ \times 100$.}
\begin{tabular}{lcccc}
\toprule
Model & 0° & 135° & 180° & NormStab$\downarrow$ \\ \midrule
ResNet18          & \textbf{86.93} & 21.28 & 36.81 & 75.67 \\
StandardCNN       & 86.09 & 24.59 & 45.58 & 71.44 \\
TernaryStdCNN     & 82.27 & 23.84 & 42.79 & 71.02 \\
\tritnet{} (full) & 80.08 & 22.58 & 41.02 & \textbf{71.80} \\
Trit-NoPredLoss   & 82.28 & ---   & ---   & 70.53 \\
Trit-NoAttention  & 81.17 & ---   & ---   & 69.69 \\
Trit-NoMemoryGate & 78.21 & ---   & ---   & 72.10 \\
\bottomrule
\end{tabular}
\end{table}

\begin{table}[H]
\centering\small
\caption{Accuracy (\%) and normalized stability on STL-10 under rotation.}
\begin{tabular}{lcccc}
\toprule
Model & 0° & 135° & 180° & NormStab$\downarrow$ \\ \midrule
ResNet18          & 55.84 & 13.51 & 30.85 & 75.80 \\
StandardCNN       & \textbf{66.86} & 16.09 & 45.59 & 75.29 \\
TernaryStdCNN     & 62.79 & 14.88 & 40.88 & 76.09 \\
\tritnet{} (full) & 60.12 & 15.76 & 37.79 & \textbf{73.78} \\
Trit-NoPredLoss   & 60.72 & ---   & ---   & 75.32 \\
Trit-NoAttention  & 61.21 & ---   & ---   & 75.52 \\
Trit-NoMemoryGate & 59.85 & ---   & ---   & 74.52 \\
\bottomrule
\end{tabular}
\end{table}

\subsection{Ablation: Triadic Structure vs.\ Ternary Quantization}

The key isolation (STL-10):
\begin{itemize}
  \item StandardCNN (float32, no triadic): NormStab = 75.29\%
  \item TernaryStdCNN (ternary, no triadic): NormStab = 76.09\%
  \item \tritnet{} (ternary + triadic): NormStab = \textbf{73.78\%}
\end{itemize}

TernaryStdCNN $\approx$ StandardCNN: \textbf{ternary weights alone do not
cause rotation robustness.} \tritnet{} outperforms both: \textbf{triadic
structure is the cause.}

Component contributions (STL-10 normalized stability improvement
vs.\ removing each component):
PredLoss $+1.54$pp, Attention $+1.73$pp, MemoryGate $+0.74$pp.

\subsection{Streaming Inference Proofs}

\begin{table}[H]
\centering\small
\caption{Memory footprint: \tritlm{} (fixed) vs.\ GPT KV-cache (linear).}
\begin{tabular}{lrrc}
\toprule
Context tokens & \tritlm{} & GPT KV-cache & Ratio \\ \midrule
100     & 0.05\,KB & 150\,KB    & 3,200$\times$ \\
1,000   & 0.05\,KB & 1,500\,KB  & 32,000$\times$ \\
10,000  & 0.05\,KB & 15,000\,KB & 320,000$\times$ \\
100,000 & 0.05\,KB & 150,000\,KB & $3.2\times10^6$$\times$ \\
\bottomrule
\end{tabular}
\end{table}

\begin{table}[H]
\centering\small
\caption{Noise recovery: loss reduction from clean trit memory
before corrupted input. Recovery increases with corruption severity.}
\begin{tabular}{lrrr}
\toprule
Corruption & No memory & Trit memory & Recovery$\uparrow$ \\ \midrule
0\%  & 2.872 & 2.824 & $+$0.049 \\
10\% & 2.951 & 2.899 & $+$0.051 \\
20\% & 3.425 & 3.409 & $+$0.016 \\
30\% & 4.170 & 3.883 & $+$0.287 \\
50\% & 7.575 & 7.231 & $+$0.345 \\
\bottomrule
\end{tabular}
\end{table}

\begin{table}[H]
\centering\small
\caption{Trit memory interpretability: cosine similarity between
memory states after reading different passage types.
Stage directions are nearly orthogonal to all other content (sim $< 0.2$).}
\begin{tabular}{lcccc}
\toprule
 & Dialogue & Stage dir. & Numbers & Punct. \\ \midrule
Dialogue    & 1.00 & 0.14 & 0.92 & 0.77 \\
Stage dir.  & 0.14 & 1.00 & 0.09 & 0.17 \\
Numbers     & 0.92 & 0.09 & 1.00 & 0.74 \\
Punct.      & 0.77 & 0.17 & 0.74 & 1.00 \\
\bottomrule
\end{tabular}
\end{table}

Forgetting resistance: after streaming 5,000 tokens of new content,
\tritlm{} performance on prior content changed by $\Delta{=}{-}0.004$
nats (slight improvement). GPT with no persistent state is unchanged
by definition; with KV-cache eviction it would degrade.

\subsection{Hardware Efficiency}

\begin{table}[H]
\centering\small
\caption{Hardware efficiency comparison.}
\begin{tabular}{lcc}
\toprule
Metric & ResNet18 & \tritnet{} \\ \midrule
Parameters      & 11,181,642 & \textbf{396,174} \\
Float32 size    & 44.7\,MB    & 1.58\,MB \\
Ternary size    & ---         & \textbf{0.078\,MB} \\
Compression     & ---         & \textbf{20.2$\times$} \\
Zero weights    & ---         & ${\sim}50\%$ \\
Energy saving   & ---         & ${\sim}87\%$~\cite{andri2018} \\
\bottomrule
\end{tabular}
\end{table}

\begin{table}[H]
\centering\small
\caption{\tritlm{} vs.\ standard LSTM per memory update step.}
\begin{tabular}{lcc}
\toprule
Metric & LSTM (float32) & TritMemoryCell \\ \midrule
Gates         & 4              & 3 \\
Weight type   & float32        & $\{-1,0,+1\}$ \\
Ops/step      & 131,072 MACs   & 384 ternary ops \\
Speedup       & ---            & ${\sim}341\times$ \\
State size    & float vector   & 25.4\,bytes \\
Interpretable & no             & yes \\
\bottomrule
\end{tabular}
\end{table}

\section{Hardware Specification}

\subsection{Primitives (36/36 tests pass)}

Trit encoding: $00=-1$, $01=0$, $10=+1$.
Consensus gate (27 cases), ternary adder with carry (6 cases),
ternary register (3 cases) --- all verified in simulation.

\subsection{Ternary MAC}

\begin{equation*}
\text{MAC}(x,w) = \sum_i
\begin{cases} +x_i & w_i=+1 \\ -x_i & w_i=-1 \\ 0 & w_i=0 \end{cases}
\end{equation*}

Zero weights are clock-gated. With ${\sim}50\%$ zero weights,
half of all MAC units are idle per inference step.

\subsection{FPGA Estimate (Xilinx Ultrascale+)}

\begin{table}[H]
\centering\small
\begin{tabular}{lrrr}
\toprule
Resource   & Estimate & Available & Util. \\ \midrule
LUTs       & ${\sim}42$k  & 600k  & 7\%   \\
Flip-Flops & ${\sim}120$k & 1.2M  & 10\%  \\
BRAM       & ${\sim}78$\,KB & 32\,MB & $<$1\% \\
Frequency  & ${\sim}250$\,MHz & --- & ---  \\
Power      & ${\sim}0.8$\,W  & --- & ---  \\
\bottomrule
\end{tabular}
\caption{Estimated FPGA resources (Vivado synthesis pending).
Compare: GPU inference ${\sim}$300\,W.}
\end{table}

\section{Discussion}

\textbf{Rotation robustness mechanism.}
The Observer stream ($1\times1$ kernel) operates at the level of individual
pixels, dynamically gating whether each location contributes via local
($3\times3$) or relational ($5\times5$) features. Under rotation, absolute
spatial positions shift but relational structure between parts is more
preserved. The triadic block learns to weight the relational stream more
heavily in ambiguous orientations, without any rotation-specific supervision.

\textbf{Memory as error correction.}
The noise recovery result (Table~4) shows an unexpected property: recovery
increases with corruption severity (+0.34 nats at 50\% vs.\ +0.05 at 0\%).
At low corruption, the noisy input is close enough to the clean input that
memory provides little additional signal. At high corruption, the input is
nearly random and memory carries most of the semantic content. This suggests
trit memory functions as a prior over expected content --- closer to Bayesian
inference than to standard RNN state.

\textbf{Interpretability.}
Stage directions (bracketed labels, imperative verbs, no narrative) produce
a memory state nearly orthogonal to all other text types
(sim $< 0.2$). Dialogue and narrative text cluster together (sim $> 0.74$).
This structure emerges from training on next-character prediction alone ---
the trit cells learn content categories without supervision. In float32
models the hidden state is a continuous vector with no discrete
interpretation; individual dimensions have no stable meaning across inputs.

\textbf{Speed on GPU vs.\ hardware.}
The sequential TritMemoryCell update is slower than GPT on GPU because
GPUs are optimized for parallel float32 operations. The theoretical
$341\times$ operation reduction over LSTM requires ternary silicon where
add costs ${\sim}$1.1\,pJ vs.\ multiply at ${\sim}$3.7\,pJ~\cite{andri2018}.
This is a hardware claim, not a software simulation claim.

\textbf{Limitations.}
Peak accuracy at 0° trails ResNet18 by 6.85pp on CIFAR-10 (86.93\% vs.\
80.08\%). The gap is attributable to capacity: \tritnet{} has 28$\times$
fewer parameters. Wider hidden channels close this gap at the cost of
the parameter efficiency advantage.

\section{Conclusion}

012 demonstrates five properties unavailable in standard binary neural
networks: rotation robustness from triadic structure (not quantization),
constant-size streaming memory, noise correction via trit state,
forgetting resistance, and discrete interpretable memory cells.
At 0.078\,MB and an estimated 0.8\,W on FPGA, \tritnet{} targets deployment
contexts binary silicon cannot reach. The full stack --- software,
hardware RTL, and all benchmarks --- is available at:
\url{https://github.com/[your-handle]/012-ternary}

\bibliographystyle{plain}
\begin{thebibliography}{99}

\bibitem{rastegari2016}
Rastegari et al.\ (2016). XNOR-Net: ImageNet classification using binary
convolutional neural networks. \textit{ECCV 2016}.

\bibitem{courbariaux2016}
Courbariaux et al.\ (2016). Binarized neural networks.
\textit{NeurIPS 2016}.

\bibitem{andri2018}
Andri et al.\ (2018). YodaNN: An architecture for ultralow power
binary-weight CNN acceleration. \textit{IEEE TCASI}.

\bibitem{rao1999}
Rao \& Ballard (1999). Predictive coding in the visual cortex.
\textit{Nature Neuroscience, 2}(1), 79--87.

\bibitem{friston2005}
Friston (2005). A theory of cortical responses.
\textit{Phil.\ Trans.\ R.\ Soc.\ B, 360}, 815--836.

\bibitem{hayes2001}
Hayes (2001). Third base. \textit{American Scientist, 89}(6), 490--494.

\bibitem{thompson2021}
Thompson et al.\ (2021). The computational limits of deep learning.
\textit{arXiv:2007.05558}.

\bibitem{he2016}
He et al.\ (2016). Deep residual learning for image recognition.
\textit{CVPR 2016}.

\bibitem{bengio2013}
Bengio et al.\ (2013). Estimating or propagating gradients through
stochastic neurons. \textit{arXiv:1308.3432}.

\bibitem{shukla2019}
Shukla et al.\ (2019). Ternary content addressable memory.
\textit{Nano Letters}.

\bibitem{ma2024}
Ma et al.\ (2024). The Era of 1-bit LLMs: All large language models
are in 1.58 bits. \textit{arXiv:2402.17764}.

\bibitem{mamba2023}
Gu \& Dao (2023). Mamba: Linear-time sequence modeling with
selective state spaces. \textit{arXiv:2312.00752}.

\end{thebibliography}

\end{document}
